% !TEX root = ../../main.tex

\section{Implementation Details}

\subsection{Software implementation}
\label{sec:software-implementation}

The manuscript documents the \texttt{tbs\_clustering\_analysis} software version
\texttt{0.1.0}, which is the version identifier currently exposed by the
repository. The reference implementation is written in Python and organizes the
workflow into tree construction, node-level distribution summaries,
child-parent and sibling testing, multiple-testing correction, and top-down
traversal. Version-specific engineering choices are recorded here rather than
in the main method narrative. In the software described here, the sibling basis
starts with the parent PCA directions and uses orthonormal padding if the
requested dimension exceeds the available parent rank; internal APIs can expose
child PCA projections, but the default basis constructor does not use them.

\subsection{Default analysis settings}

\Cref{tab:defaults} summarizes the default settings used in the software
described here. These settings define how the tree is built, how evidence is
combined across the tree, and how final cluster decisions are made. The table
is descriptive rather than justificatory; \Cref{tab:method-constants} records
which defaults are analytically motivated, which are empirical working choices,
and which still require validation.

\begin{longtable}{@{}p{0.36\linewidth}p{0.58\linewidth}@{}}
\caption{Default analysis settings used in the software documented here.}
\label{tab:defaults}\\
\toprule
Component & Current choice \\
\midrule
\endfirsthead
\toprule
Component & Current choice \\
\midrule
\endhead
\bottomrule
\endfoot
Tree distance & Hamming by default; explicit continuous benchmark cases use precomputed Euclidean distances \\
Linkage & Average linkage \\
Edge significance level & $\alpha_{\mathrm{edge}} = 0.001$ \\
Sibling significance level & $\alpha_{\mathrm{sib}} = 0.01$ \\
Edge test & Projected Wald test using the parent-specific orthonormal PCA basis \\
Edge dimension rule & Marchenko--Pastur signal count from the parent's local spectrum \\
Sibling test & Projected Wald test; empirical calibration is active with internal support and closes focal sibling gates when support is missing; selected-tail validity remains unproven \\
Sibling dimension rule & Geometric mean of positive child edge dimensions \\
Sibling basis / inflation & Parent PCA head capped to available rank; projected chi-square raw test; empirical inflation $\hat c$ adjusts the plug-in chi-square tail; selected-tail validity remains open \\
Minimum spectral dimension & 2 \\
Spectral rows & Descendant leaves only \\
Edge multiple-testing correction & Hierarchical BH correction on the tree \\
Sibling multiple-testing correction & Flat BH on focal pairs \\
Pass-through traversal & On \\
\end{longtable}

Two aspects of these defaults deserve emphasis. First, the two tests do not
combine information in the same way. The child-parent edge test asks whether a
child leaves the background defined by its parent. The sibling test then asks
whether the two children are different enough from one another to justify
keeping the split. Second, once those test tables exist, the final clustering
step is direct: the method walks down the tree and keeps only the splits that
remain supported after correction.

The software follows the representation contract defined in the method:
Bernoulli coordinates enter directly, categorical variables enter as one-hot
blocks with drop-last multinomial covariance, and explicit continuous inputs
enter as empirical-Gaussian blocks. Discretized Gaussian representations are
benchmark variants only.
